\documentclass[12pt]{article}

% a template that a friend gave, it's worked well enough for me
% i have added some packages and stuff that have proved useful

\usepackage{fancyhdr}
\usepackage{tipa}
\usepackage{fontspec}
\usepackage{amsfonts}
\usepackage{enumitem}
\usepackage[margin=1in]{geometry}
\usepackage{graphicx}
\usepackage{float}
\usepackage{amsmath}
\usepackage{braket}
\usepackage{amssymb}
\usepackage{booktabs}
\usepackage{hyperref}
\usepackage{mathtools}
\usepackage{xcolor}
\usepackage{float}
\usepackage{algpseudocodex}
\usepackage{titlesec}
\usepackage{bbm}

\pagestyle{fancy}
\fancyhf{} % sets both header and footer to nothing
\lhead{Kevin Sheng}
\setmainfont{Comic Neue}
\renewcommand{\headrulewidth}{1pt}
\setlength{\headheight}{0.75in}
\setlength{\oddsidemargin}{0in}
\setlength{\evensidemargin}{0in}
\setlength{\voffset}{-.5in}
\setlength{\headsep}{10pt}
\setlength{\textwidth}{6.5in}
\setlength{\headwidth}{6.5in}
\setlength{\textheight}{8in}
\renewcommand{\headrulewidth}{0.5pt}
\renewcommand{\footrulewidth}{0.3pt}
\setlength{\textwidth}{6.5in}
\usepackage{setspace}
\usepackage{multicol}
\usepackage{float}
\setlength{\columnsep}{1cm}
\setlength\parindent{24pt}
\usepackage [english]{babel}
\usepackage [autostyle, english = american]{csquotes}
\MakeOuterQuote{"}

\setlength{\parskip}{6pt}
\setlength{\parindent}{0pt}

\titlespacing\section{0pt}{12pt plus 4pt minus 2pt}{0pt plus 2pt minus 2pt}
\titlespacing\subsection{0pt}{12pt plus 4pt minus 2pt}{0pt plus 2pt minus 2pt}
\titlespacing\subsubsection{0pt}{12pt plus 4pt minus 2pt}{0pt plus 2pt minus 2pt}

\hypersetup{colorlinks=true, urlcolor=blue}

\newcommand{\correction}[1]{\textcolor{red}{#1}}


\rhead{Math 110A}

\makeatletter
\def\@seccntformat#1{%
  \expandafter\ifx\csname c@#1\endcsname\c@section\else
  \csname the#1\endcsname\quad
  \fi}
\makeatother

\newcommand{\lra}{\xLeftrightarrow}
\newcommand{\ra}{\xRightarrow}
\DeclareMathOperator{\Fun}{Fun}
\DeclareMathOperator{\id}{id}
\DeclareMathOperator{\pr}{pr}

\begin{document}

\section{Problem 1}

From class, we've shown that a homomorphism of $R$-algebras $h$
that starts from $R\braket{x_i, i \in I}$ is defined by its behavior on individual $x_i$.
If we fix $h(x_i)=x_i$, then the values of $h$ on all
other noncommutative polynomials must be pinned down.

In class, we also saw that this behavior can be written as
\[h\left(\sum_{\phi} r_\phi x^\phi\right)=\sum_{\phi} r_\phi \prod_{i=0}^{|\phi|} x_{\phi(i)}\]
where the only difference is that the order of the $x_i$s on the RHS doesn't matter.

To show that this homomorphism is surjective, fix a polynomial in $R[x_i, i \in I]$.
Even though it's commutative, we can still write it as $\sum_{\phi} r_\phi x^\phi$.
We can then just take this as our preimage, since $h$ doesn't touch coefficient
nor the behavior of any monomials.
It really just makes order not matter anymore.

If $|I| > 2$, though, this homomorphism isn't injective.
Consider $x_1x_2$ and $x_2x_1$, different elements under $R\braket{x_i, i \in I}$.
If we put them through $h$, both come out to the same polynomial $x_1x_2$ under
$R[x_i, i \in I]$.

\pagebreak

\section{Problem 2}

I think for this problem it's easier to define our polynomials by thinking
about the equivalence classes on $\phi$.
Call monomials $\phi_1$ and $\phi_2$ equivalent if $c: ([n] \to I) \to \N^I$, defined by
\[c(\phi)_i=\sum_{j=1}^{n} \mathbf{1}[\phi(j)=i]\]
gives the same output on both of them.
In other words, our monomials should now commute.
I'll use $\sum_{[\psi]} r_i x^\psi$ when discussing these commutative polynomials.

\subsection{Injective}

Consider two homomorphisms of $R$-algebras $h_1$ and $h_2$ s.t. $h_1(x_i)=h_2(x_i)$ for all $x_i \in I$.

Now consider $h_1$ evaluated on an arbitrary polynomial $\sum_{[\psi]} r_i x^\psi$:
\begin{align*}
  h_1\left(\sum_{[\psi]} r_i x^\psi\right)
   & = \sum_{[\psi]} \phi(r_i) \prod_{i=1}^{|\psi|} h_1(\psi(i)) \\
   & = \sum_{[\psi]} \phi(r_i) \prod_{i=1}^{|\psi|} h_2(\psi(i)) \\
   & = h_2\left(\sum_{[\psi]} r_i x^\psi\right)
\end{align*}
Since we're mapping to $S^I_{\text{comm}}$, the exact representatives
of the equivalence classes don't matter and this chain of logic is still well-defined.

\pagebreak

\subsection{Surjective}

Fix an $(s_i)_{i \in I} \in S^I_{\text{comm}}$.

I propose the following homomorphism of $R$-algebras:
\[h\left(\sum_{[\psi]} r_i x^\psi\right)
  =\sum_{[\psi]} \phi(r_i) \prod_{i=1}^{|\psi|} h(\psi(i))\]
defined off of $h(\psi(i))=(s_i)_{\psi(i)}$.
It should be clear that on individual $x_i \in I$ this mapping satisfies
$h(x_i)=\phi(1) \cdot (s_i)_{x_i} = (s_i)_{x_i}$.

Now we NTS this is actually, well, a homomorphism of $R$-algebras.
\begin{enumerate}
  \item \textbf{Addition:}
        \begin{align*}
          h\left(\sum_{[\psi]} r_i x^\psi+\sum_{[\psi]} s_i x^\psi\right)
           & = h\left(\sum_{[\psi]} (r_i+s_i) x^\psi\right)                       \\
           & =\sum_{[\psi]} \phi(r_i+s_i) \prod_{i=1}^{|\psi|} h(\psi(i))         \\
           & =\sum_{[\psi]} (\phi(r_i)+\phi(s_i)) \prod_{i=1}^{|\psi|} h(\psi(i)) \\
           & = \sum_{[\psi]} \phi(r_i) \prod_{i=1}^{|\psi|} h(\psi(i))
          + \sum_{[\psi]} \phi(s_i) \prod_{i=1}^{|\psi|} h(\psi(i))               \\
           & = h\left(\sum_{[\psi]} r_i x^\psi\right)
          +h\left(\sum_{[\psi]} s_i x^\psi\right)
        \end{align*}

  \item \textbf{Multiplication:}
        \begin{align*}
          h\left(\sum_{[\psi]} r_i x^\psi \cdot \sum_{[\psi]} s_i x^\psi\right)
           & = h\left(\sum_{[\psi_1]} \sum_{[\psi_2]} r_i s_i x^{\psi_1 \sqcup \psi_2}\right)                                                    \\
           & = \sum_{[\psi_1]} \sum_{[\psi_2]} \phi(r_i s_i) h\left(x^{\psi_1 \sqcup \psi_2}\right)                                              \\
           & = \sum_{[\psi_1]} \sum_{[\psi_2]} \phi(r_i)\phi(s_i) \prod_{i=1}^{|\psi_1|+|\psi_2|} h((\psi_1 \sqcup \psi_2)(i))                   \\
           & = \sum_{[\psi_1]} \sum_{[\psi_2]} \phi(r_i)\phi(s_i) \prod_{i=1}^{|\psi_1|} h(\psi_1(i)) \prod_{i=1}^{|\psi_2|} h(\psi_2(i))        \\
           & = \sum_{[\psi_1]} \phi(r_i) \prod_{i=1}^{|\psi_1|} h(\psi_1(i)) \cdot \sum_{[\psi_2]} \phi(r_i) \prod_{i=1}^{|\psi_2|} h(\psi_2(i)) \\
           & = h\left(\sum_{[\psi]} r_i x^\psi\right) \cdot h\left(\sum_{[\psi]} s_i x^\psi\right)
        \end{align*}

  \item \textbf{One:} \\
        Since $\phi$ itself is a homomorphism, $h(1)=\phi(1)=1$.

  \item \textbf{Preservation:} \\
        If we take any $r \in R$, we indeed have $h(r)=\phi(r)$
        since $r$ doesn't have any coefficients.
\end{enumerate}

\section{Problem 3}

Let $s_1, s_2 \in S$.
Since $f$ is surjective, $\exists t_1, t_2 \in T: f(t_1)=s_1, f(t_2)=s_2$.

Then,
\begin{align*}
  s_1s_2
   & = h(t_1)h(t_2)        \\
   & = h(t_1 \cdot t_2)    \\
   & = h(t_2 \cdot t_1)    \\
   & = h(t_2) \cdot h(t_1) \\
   & = s_2s_1\quad\square
\end{align*}

\section{Problem 4}

We at least need $2$ matrices to generate a sort of spanning set with multiplication.

Let's have our matrices be
\begin{align*}
  A=\begin{bmatrix}
      0 & 1 \\
      1 & 0
    \end{bmatrix} &  &
  B=\begin{bmatrix}
      1 & 0 \\
      0 & 0
    \end{bmatrix}
\end{align*}
Then, we can multiply them to get these two matrices:
\begin{align*}
  AB=\begin{bmatrix}
       0 & 0 \\
       1 & 0
     \end{bmatrix} &  &
  ABA=\begin{bmatrix}
        0 & 0 \\
        0 & 1
      \end{bmatrix}
\end{align*}
and all four of these combined can form a basis for all 2 by 2 matrices.

To show that this is the minimum, BWOC say we could generate all 2 by 2 matrices
by just raising a single matrix $A$ to various powers.
In that case, we could express any matrix as $\sum_{i=0}^{n} \lambda_i A^i$.

Taking two matrices $M_1$ and $M_2$ written in this manner, we have
\begin{align*}
  M_1M_2
   & = \sum_{i=0}^{n} \lambda_i A^i \cdot \sum_{i=0}^{n} \lambda'_i A^i \\
   & = \sum_{i=0}^{2n} \sum_{j=0}^{i} \lambda_j \lambda'_{i-j} A^i      \\
   & = \sum_{i=0}^{2n} \sum_{j=0}^{i} \lambda'_j \lambda_{i-j} A^i      \\
   & = \sum_{i=0}^{n} \lambda'_i A^i \cdot \sum_{i=0}^{n} \lambda_i A^i \\
   & = M_2M_1
\end{align*}
which would then mean that all matrices commute with each other, a contradiction.

\section{Problem 5}

one must imagine sisyphus happpy

\subsection{Addition Commutes}

We have this by commutativity in each of the $R_i$s:
\[(r_i)_{i \in I}+(r'_i)_{i \in I}=(r_i+r'_i)_{i \in I}=(r'_i+r_i)_{i \in I}=(r'_i)_{i \in I}+(r_i)_{i \in I}\]

\subsection{Addition Associates}

By associativity on the individual rings,
\begin{align*}
  (r_i)_{i \in I}+((s_i)_{i \in I}+(t_i)_{i \in I})
   & = (r_i+(s_i+t_i))_{i \in I}                         \\
   & = ((r_i+s_i)+t_i)_{i \in I}                         \\
   & = ((r_i)_{i \in I}+(s_i)_{i \in I})+(t_i)_{i \in I}
\end{align*}
The proof that multiplication is associative goes basically the exact same way,
save for using $\cdot$ instead of $+$.

\subsection{Addition is a Bijection}

Fix an $(r_i)_{i \in I}$ and consider addition with it as a function.

To prove injectivity, take $(s_i)_{i \in I}$ and $(t_i)_{i \in I}$ s.t. $(r_i)_{i \in I}+(s_i)_{i \in I}=(r_i)_{i \in I}+(t_i)_{i \in I}$.
Then, since addition within each individual ring makes a bijection,
\[(r_i+s_i)_{i \in I}=(r_i+t_i)_{i \in I}
  \implies (s_i)_{i \in I}=(t_i)_{i \in I}\]

For surjectivity, fix a $(s_i)_{i \in I}$ that we want to hit.
Then,
\begin{align*}
  (r_i)_{i \in I}+(s_i-r_i)_{i \in I}
   & = (r_i+s_i-r_i)_{i \in I} \\
   & = (s_i)_{i \in I}
\end{align*}

\subsection{Distributivity}

By distributivity in each of the $R_i$s, left distributivity follows:
\begin{align*}
  (r_i)_{i \in I} \cdot ((s_i)_{i \in I}+(t_i)_{i \in I})
   & = (r_i \cdot (s_i+t_i))_{i \in I}                                         \\
   & = (r_i \cdot s_i+r_i \cdot t_i)_{i \in I}                                 \\
   & = (r_i \cdot s_i)_{i \in I}+(r_i \cdot t_i)_{i \in I}                     \\
   & = (r_i)_{i \in I}\cdot(s_i)_{i \in I}+(r_i)_{i \in I}\cdot(t_i)_{i \in I}
\end{align*}
and so does right distributivity:
\begin{align*}
  ((r_i)_{i \in I}+(s_i)_{i \in I})\cdot(t_i)_{i \in I}
   & = ((r_i+s_i)\cdot t_i)_{i \in I}                                          \\
   & = (r_i\cdot t_i+s_i\cdot t_i)_{i \in I}                                   \\
   & = (r_i\cdot t_i)_{i \in I}+(s_i\cdot t_i)_{i \in I}                       \\
   & = (r_i)_{i \in I}\cdot(t_i)_{i \in I}+(s_i)_{i \in I}\cdot(t_i)_{i \in I}
\end{align*}

\subsection{Identities}

Our zero is $(0_i)_{i \in I}$, where $0_i$ is the additive identity in $R_i$:
\begin{align*}
  (r_i)_{i \in I}+(0_i)_{i \in I}
   & = (r_i+0_i)_{i \in I} \\
   & = (r_i)_{i \in I}
\end{align*}
Similarly, our $1$ is $(1_i)_{i \in I}$, where $1_i$ is the additive identity in $R_i$:
\begin{align*}
  (r_i)_{i \in I} \cdot (1_i)_{i \in I}
   & = (r_i \cdot 1_i)_{i \in I} \\
   & = (r_i)_{i \in I}
\end{align*}

\subsection{Projection is a Homomorphism}

Fix a $j \in I$.

First off, $\pr_j((1_i)_{i \in I})=1_j$, so one still maps to one.

For addition,
\begin{align*}
  \pr_j((r_i)_{i \in I}+(s_i)_{i \in I})
   & =\pr_j((r_i+s_i)_{i \in I})                     \\
   & = r_j+s_j                                       \\
   & = \pr_j((r_i)_{i \in I})+\pr_j((s_i)_{i \in I})
\end{align*}
and for multiplication,
\begin{align*}
  \pr_j((r_i)_{i \in I} \cdot (s_i)_{i \in I})
   & = \pr_j((r_i \cdot s_i)_{i \in I})                    \\
   & = r_j \cdot s_j                                       \\
   & = \pr_j((r_i)_{i \in I}) \cdot \pr_j((s_i)_{i \in I})
\end{align*}

\subsection{Commutativity Conditions}

If all $R_i$s are commutative,
\begin{align*}
  (r_i)_{i \in I} \cdot (s_i)_{i \in I}
   & = (r_i \cdot s_i)_{i \in I}             \\
   & = (s_i \cdot r_i)_{i \in I}             \\
   & = (s_i)_{i \in I} \cdot (r_i)_{i \in I}
\end{align*}

OTOH, if they weren't, $\exists j \in I$ and $r, s \in R_j$ s.t. $r \cdot s \ne s \cdot r$.

Then we can let $(r_i)_{i \in I}$ be any $I$-tuple s.t. $r_j=r$ and $(s_i)_{i \in I}$ be defined similarly.
At index $j$, $(r_i)_{i \in I} \cdot (s_i)_{i \in I}$ is $rs$, while
$(s_i)_{i \in I} \cdot (r_i)_{i \in I}$ is $sr$.
This makes the two products different and $\prod_{i \in I} R_i$ not commutative.

\section{Problem 6}

Our isomorphism $h$ should be defined by $(r_x)_{x \in X} \to f(x)=r_x$.

First let me show that this is a homomorphism though:
\begin{enumerate}
  \item For $(1)_{x \in X}$, $f(x)=1$ for all $x$, which makes it the one of $\Fun_R(X)$.

  \item $(r_x)_{x \in X}+(r'_x)_{x \in X}$ maps to $f(x)=r_x+r'_x$.
        The two elements individually map to $f(x)=r_x$ and $f(x)=r'_x$,
        two functions which add up pointwise to our original $f(x)$.

  \item $(r_x)_{x \in X} \cdot (r'_x)_{x \in X}$ maps to $f(x)=r_x \cdot r'_x$.
        The two elements individually map to $f(x)=r_x$ and $f(x)=r'_x$,
        two functions which when multiplied pointwise become our original $f(x)$.

  \item The homomorphism from $R$ to $\Fun_X(R)$ takes $r$ to $f(x)=r$ for all $x$.
        When taken to $\prod_{x \in X} R$ first, it becomes $(r)_{x \in X}$.
        After applying our homomorphism, it once again becomes $f(x)=r$.
\end{enumerate}

To upgrade it to an isomorphism, we NTS that it's a bijection.

For injectivity, notice that if $(r_x)_{x \in X} \ne (r'_x)_{x \in X}$,
$\exists i \in X: r_i \ne r'_i$.
Then $h((r_x)_{x \in X})(i)=r_i$ and $h((r'_x)_{x \in X})(i)=r'_i$,
so $h((r_x)_{x \in X}) \ne h((r'_x)_{x \in X})$.

For surjectivity, fix an $f: X \to R$.
If we define $r_x=f(x)$ for each $x$, then $h((r_x)_{x \in X})(x)=f(x)$ and $h((r_x)_{x \in X})=f$.

\end{document}
